% Transcription — fonds Alexandre Grothendieck, shelfmark 160, batch 1 (pages 1-20).
% Established from the facsimile archives/batches/160/batch-01.pdf (a mirror of
% https://grothendieck.umontpellier.fr/160.pdf), per the skill
% transcribe-grothendieck. No cover sheet: PDF page N = archivists' page N,
% checked against the pencilled numbers bottom-left.
% Page 17 is blank and carries no mathematics — absent.
% Pass: Opus 5 (claude-opus-5), 2026-09-19 — first pass, unchecked against the pages by a human.
\documentclass[11pt,a4paper]{article}
% Le préambule commun à toutes les transcriptions du fonds Grothendieck.
%
% Il tient en une idée : **l'incertitude se compose, elle ne se résout pas.**
% Une transcription qui lisse un mot douteux a détruit la seule chose pour
% laquelle on la faisait — un lecteur doit pouvoir savoir, à chaque endroit,
% ce qui était sur le feuillet et ce qui a été ajouté. Les six macros qui
% suivent sont donc l'appareil critique, et non de la décoration.
%
% Les mêmes commandes sont reconnues par `scripts/render.mjs`, qui produit la
% vue de lecture du site. Une macro ajoutée ici doit l'être aussi là-bas,
% faute de quoi le rendu échouera bruyamment — ce qui est voulu : un rendu
% silencieusement faux est pire qu'un rendu absent.

\ProvidesPackage{grothendieck}[2026/08/08 Transcriptions du fonds Grothendieck]

\usepackage{amsmath,amssymb,amsthm}
\usepackage{mathrsfs}
\usepackage{stmaryrd}
% Les diagrammes commutatifs sont partout dans ce fonds : c'est la langue
% même de la géométrie algébrique de Grothendieck, et une transcription qui
% les rendrait en prose ne transcrirait rien.
\usepackage{tikz-cd}
% Français par défaut — c'est la langue des feuillets — mais l'anglais est
% chargé aussi : la traduction et le résumé sont des documents anglais, et la
% typographie française (espaces avant les deux-points) y serait une faute.
% Ces documents basculent par \selectlanguage{english} en tête de corps.
\usepackage[english,french]{babel}
\usepackage{fontspec}
\usepackage{geometry}
\geometry{a4paper,margin=25mm,marginparwidth=28mm}
\usepackage{marginnote}
\usepackage{xcolor}
% ulem, not soul. `soul`'s \st and \ul are built on a character-by-character
% scan that XeLaTeX's font machinery breaks: the document fails with an
% undefined control sequence rather than degrading. ulem does the same two jobs
% and is Unicode-engine safe. `normalem` stops it hijacking \emph, which the
% transcriptions use for Grothendieck's own emphasis.
\usepackage[normalem]{ulem}
\usepackage{hyperref}
\hypersetup{colorlinks=true,linkcolor=black,urlcolor=black,pdfborder={0 0 0}}

% --- Métadonnées du lot -----------------------------------------------------
% Reprises telles quelles par le script de rendu, qui les affiche en tête de
% la vue de lecture. Elles fixent ce que le fichier prétend couvrir : sans
% elles, une transcription flotte sans point d'attache dans un fonds de seize
% mille feuillets.

\newcommand{\folder}[1]{\gdef\grothFolder{#1}}
\newcommand{\batch}[1]{\gdef\grothBatch{#1}}
\newcommand{\leaves}[2]{\gdef\grothFirst{#1}\gdef\grothLast{#2}}
\newcommand{\foldertitle}[1]{\gdef\grothFoldertitle{#1}}
\gdef\grothFolder{?}\gdef\grothBatch{?}\gdef\grothFirst{?}\gdef\grothLast{?}\gdef\grothFoldertitle{}

% --- Le résumé --------------------------------------------------------------
%
% Il ouvre la lecture modernisée, sous le titre « Résumé », et il est composé
% en petit. Ce n'est pas un ornement : c'est le seul passage du document qui
% ne soit pas des mathématiques, et un lecteur doit pouvoir le reconnaître
% avant de l'avoir lu — puis le sauter s'il connaît déjà le sujet, ou s'y
% arrêter s'il ne le connaît pas. Le filet qui le suit marque la reprise du
% corps.

\newenvironment{resume}
  {\small\setlength{\parskip}{0.45em}}
  {\par\medskip\hrule\medskip}

% --- Mots-clés ---------------------------------------------------------------
%
% En anglais, en clôture du résumé de la lecture modernisée : le vocabulaire
% moderne sous lequel chercher ce que les feuillets construisent. Le manifeste
% du site (`npm run manifest`) les extrait pour étiqueter la cote entière —
% cette ligne est la source unique des tags, il n'y a pas de fichier à côté.
\newcommand{\keywords}[1]{\par\smallskip\noindent
  {\footnotesize\textsc{Keywords} --- \textit{#1}}\par}

% --- L'appareil critique ----------------------------------------------------

% Le numéro de feuillet, en marge. C'est l'ancre qui permet de régler un
% désaccord sur une formule en regardant *un* feuillet plutôt qu'une chemise
% de six cents. Le site s'en sert aussi pour tourner les pages du fac-similé
% quand on fait défiler la transcription.
\newcommand{\leaf}[1]{%
  \marginnote{\footnotesize\color{gray!70}\textsf{#1}}%
  \label{leaf:#1}\ignorespaces}

% Illisible : on ne devine pas. Un mot inventé qui se lit comme les autres est
% le pire résultat possible.
\newcommand{\ill}{\textcolor{red!60!black}{[\,\ldots\,]}}

% Lecture douteuse : le mot est donné, et signalé comme incertain.
\newcommand{\uncertain}[1]{\uline{#1}}

% Ajout de l'éditeur — une lettre manquante, un mot sous-entendu.
\newcommand{\add}[1]{\textcolor{blue!50!black}{[#1]}}

% Biffé par Grothendieck lui-même. Ce qu'il a rayé fait partie du document :
% une rature dit souvent où la pensée a changé de direction.
\newcommand{\struck}[1]{\sout{#1}}

% Note du transcripteur, sur l'état matériel du feuillet.
\newcommand{\note}[1]{\par\smallskip{\small\color{gray!60!black}\textit{[#1]}}\par\smallskip}

% Note marginale de Grothendieck, distincte du corps du texte.
\newcommand{\marginal}[1]{\par\smallskip\noindent\hspace{1em}%
  {\small\color{gray!60!black}#1}\par\smallskip}

% --- Titre ------------------------------------------------------------------

\AtBeginDocument{%
  \begin{center}
    {\large\bfseries Fonds Alexandre Grothendieck --- cote n\textsuperscript{o}~\grothFolder}\\[2pt]
    {\itshape\grothFoldertitle}\\[4pt]
    {\small feuillets \grothFirst--\grothLast\ (lot \grothBatch)}\\[2pt]
    {\footnotesize\color{gray!60!black}Transcription établie d'après le fac-similé numérisé
     par l'Université de Montpellier.\\
     Les lectures incertaines sont soulignées, les illisibles notées [\,\ldots\,],
     les ajouts entre crochets.}
  \end{center}
  \vspace{1em}}

\endinput


\folder{160}
\batch{1}
\pages{1}{20}
\dating{s.d. — le groupe « Dérivateurs » (157-1 à 160) est daté 1990-[1991]}
\watermark{Édition de démonstration}
\foldertitle{Schémas [en] groupes : notes manuscrites (s.d.)}

\begin{document}

\section*{Un lemme d'algébrisation le long d'un fermé}

\page{1}

\begin{tikzcd}
X' \arrow[d, "p"] \\ X
\end{tikzcd}

\textbf{Lemme.} Soit $f : X' \to X$ morphisme \struck{propre} \add{de} type fini
de schémas noethériens, $V'$ l'ouvert des points de $X'$ où $f$ est étale, $Z$ une
partie fermée de $X$, $\hat{X}$ le complété formel de $X$ le long de $Z$,
$q$ (ou $(q_n)$) une section \uncertain{formelle} de $X'$ sur $X$ \ill{}
(i.e. une collection de sections des $X'\times_X X_n$ sur les $X_n$),
$Z' = q_\bullet(\hat{X}_\bullet)$, \uncertain{$= q(\ldots)$}
(qui est une sous-section \ill{} de $X'$ \ill{}, avec $p|Z' : Z' \simeq Z$).

1) On peut \ill{}
\[
  U_Z = q^{-1}(V') \longleftarrow V' \cap Z' .
\]
\note{la flèche est tracée vers la gauche sur le feuillet}

a) $X_n \cap U_Z = q_n^{-1}(V')$ \uncertain{est schématiquement dense dans $\hat{X}$}

b) $X'$ est l'\uncertain{adhérence schématique} \ill{} de la famille des
$q_n(X_n) \subset X'$, \ill{}

c) La famille des $V'_n$ \add{(}\uncertain{est schématiquement dense dans $X'$}\add{)}
\marginal{i.e. en un $\hat{V'}$}

2) Cas \uncertain{où} $V' \supset Z$, i.e. $f$ est étale aux points de $Z$
\supplied{(}et il existe une partie ouverte $U' \subset V'$ \uncertain{contenant} $Z$,
\ill{} pour \uncertain{que} $U'\times_X Z \xrightarrow{\ \sim\ } Z' \xrightarrow{\ \sim\ } Z$.

\textbf{Dém.} \struck{OPS $X$ affine, \ill{} hypothèse \ill{} par localisation}
La condition c) équivaut : \ill{} des conditions

$c_1$) $V'$ schématiquement dense dans $X'$,

$c_2$) $\hat{V'}$ schématiquement dense dans $U'$, i.e.\ill{}

\section*{Le schéma des sous-groupes finis de rang $p^n$ de $\mathbb{G}_a$}

\page{2}
$\mathbb{F}_p$, $\mathbb{G}_a = \mathbb{G}_{a\,\mathbb{F}_p}$, $n \in \mathbb{N}$,
$X_n$ schéma des sous-groupes finis de rang $p^n$ de $\mathbb{G}_a$,
\[
  \mathbb{E}^n = \mathbb{E}^n_{\mathbb{F}_p} \xrightarrow{\ \varphi_n\ } X_n
\]
\[
  a = (a_0, \dots, a_{n-1}) \longmapsto \operatorname{Ker}(f_a)
\]
\[
  f_a : \mathbb{G}_a \longrightarrow \mathbb{G}_a, \qquad
  f_a(x) = x^{p^n} + a_{n-1}x^{p^{n-1}} + \dots + a_1 x^p + a_0 x .
\]

\textbf{Lemme.} $\varphi_n$ est une \struck{\ill{}} \add{immersion fermée}
surjective, donc induit $\mathbb{E}^n \xrightarrow{\ \sim\ } X_{n\,\mathrm{réd}}$.
\note{le correctif « immersion fermée » est porté au-dessus d'un mot biffé}

a) \emph{Surjectif} : il suffit de prouver que
$\mathbb{E}^n \to X_n(k)$ surj. si $k$ alg. clos.
Or si $H \subset \mathbb{G}_{a,k}$, $H \in X_n(k)$ i.e. $H$ fini de rang $p^n$, \ill{}
\[
  H = \alpha_p^r \times (\mathbb{Z}/p)^s_k, \qquad r+s = n .
\]
Soit $\xi_1, \dots, \xi_s$ base de $H(k)$ sur $\mathbb{F}_p$.
\marginal{$H^0$, $H(k)_k$}
On a $H(k)_k = \operatorname{Ker}(f_\alpha)$, $\alpha \in \mathbb{E}^s_k$, \uncertain{à justifier} :
\[
  \alpha = (\alpha_0, \dots, \alpha_{s-1}), \qquad
  f_\alpha(x) = x^{p^s} + \alpha_{s-1}x^{p^{s-1}} + \dots + \alpha_1 x^{p} + \alpha_0 x .
\]
Considérons alors
\[
  f : \mathbb{G}_a \longrightarrow \mathbb{G}_a, \qquad f = f_\alpha \circ \mathrm{Frob}^r
\]
\[
  f(x) = x^{p^{r+s}} + \alpha_{s-1} x^{p^{r+s-1}} + \dots + \alpha_0 x^{p^r}
\]
\note{il annote les coefficients : $a_{r+s-1}$, $a_r$, et $a_{r-1} = \dots = a_0 = 0$}
On a
\[
  \operatorname{Ker} f = f_\alpha^{-1}(\alpha_{p^r}) \supset H(k)_k \cdot \alpha_{p^r}
\]
\note{sous $\alpha_{p^r}$ il écrit « $\operatorname{Ker}\mathrm{Frob}^r$ »}
et comme les $2$ \uncertain{membres} sont de \ill{} rang $p^{r+s} = p^n$,
\uncertain{l'inclusion} $\supset$ est une égalité.

b) \textbf{Mono.} Soient \struck{\ill{}} $k$ $\mathbb{F}_p$-alg.,
$a, b \in \mathbb{E}^n$ tels que $\varphi_n(a) = \varphi_n(b)$, i.e.
$\operatorname{Ker} f_a = \operatorname{Ker} f_b$ ; prouvons $f_a = f_b$ (i.e. $a = b$)
\struck{\ill{}}

\begin{tikzcd}[column sep=small, row sep=small, nodes={font=\scriptsize}]
H \arrow[r, hook] & \mathbb{G}_{a,k} \arrow[r, bend left=20, "f_a"] \arrow[r, bend right=20, "f_b"'] \arrow[rd, "g_a"'] & \mathbb{G}_{a,k} \\
 & & \mathbb{G}_{a,k}/H
\end{tikzcd}

\note{la seconde flèche descendante porte les deux noms $g_a$, $g_b$}

$g_b = u\,g_a$ donc $f_b = u f_a$, $u \in \operatorname{Aut}(\mathbb{G}_{a,k})$.
Il faut prouver $u = \operatorname{id}_{\mathbb{G}_a}$. Clair si $k$ réduit, car
\uncertain{$\operatorname{Aut}(\mathbb{G}_a) = k^*$}

\page{3}
Cas général : \struck{OPS} $k$ noethérien, et \uncertain{que} $J$ idéal
\struck{\ill{}} de carré nul, $k_0 = k/J$, \struck{\ill{}} $u_0 = {}$\ill{}

\textbf{Lemme.} Soit $a \in \mathbb{E}^n$, $u \in \operatorname{Aut}(\mathbb{G}_{a,k})$
tel que $u \circ f_a$ soit de la forme $f_b$ $(b \in k^n)$, alors
$u = \operatorname{id}_{\mathbb{G}_{a,k}}$.

Cas général : OPS $k$ noethérien, $J$ idéal de carré nul
et $u$ tel que $u_0 = \operatorname{id}_{\mathbb{G}_{a,k_0}}$ $(k_0 = k/J)$. Alors
\[
  u(x) = \bigl(x + c_1 x^{p} + c_2 x^{p^2} + \dots + c_N x^{p^N}\bigr), \qquad c_i \in J
\]
\note{il écrit d'abord $c_0 x$ puis ramène le coefficient de $x$ à $1$ ; voir la fin de la page}
\begin{align}
  u \circ f_a(x) &= c_0\bigl(x^{p^n} + a_{n-1}x^{p^{n-1}} + a_{n-2}x^{p^{n-2}} + \dots + a_1 x^p + a_0 x\bigr) \notag \\
   &\quad {} + c_1\bigl(x^{p^{n+1}} + a_{n-1}^{p} x^{p^{n}} + \dots + a_0^{p} x^{p}\bigr) \notag \\
   &\quad {} + c_2\bigl(x^{p^{n+2}} + a_{n-1}^{p^2} x^{p^{n+1}} + \dots + a_0^{p^2} x^{p^2}\bigr) \notag \\
   &\quad \dots \notag \\
   &\quad {} + c_N\bigl(x^{p^{n+N}} + a_{n-1}^{p^N} x^{p^{n+N-1}} + \dots + a_0^{p^N} x^{p^N}\bigr) \notag
\end{align}
Le coeff. de $x^{p^{n+N}}$ dans l'expression \ill{} est \uncertain{$c_N x^{p^{n+N}}$},
donc pour $N \geqslant 1$, on doit avoir $c_N = 0$. On en déduit de proche en proche
$c_{N-1} = 0, \dots, c_1 = 0$. Donc $u(x) = c_0(x) = c_0 x^{p^n}$ \ill{}, et il
faut aussi $c_0 = 1$. cqfd.

c) \emph{$\varphi_n$ \struck{est} propre.} \ill{} \uncertain{il faut prouver} \ill{}
\struck{\ill{}}

\textbf{Lemme.} Soit $S$ \uncertain{un trait} sur $\mathbb{F}_p$,
$H \subset \mathbb{G}_{a,S}$, $H \in X_n(S)$ i.e. $H$ plat fini $/S$ de rang $p^n$ ;
supposons que $H_\eta = \operatorname{Ker} f_{a_\eta}$, $a_\eta \in K^n$
$(K = k(\eta))$, alors $a_\eta \in V^n$ et \struck{\ill{}}
$H = \operatorname{Ker}\bigl(f_a : \mathbb{G}_{a,S} \to \mathbb{G}_{a,S}\bigr)$.
\note{$V$ est l'anneau de valuation du trait $S$}

\page{4}
Supposons que $a_\eta = (a_0, a_1, \dots, a_{n-1})$ soit $\notin V^n$,
\struck{\ill{}} soit $\nu = -\inf \operatorname{val}(a_i)$, $\pi$ uniformisante de $V$,
soit $b = \pi^\nu a = (\pi^\nu a_0, \dots, \pi^\nu a_{n-1})$, de sorte que les
$b_i = \pi^\nu a_i$ \struck{soit} $(0 \leqslant i \leqslant n-1)$ \uncertain{soient} $\in V$
et l'un est une unité. On a, si
\[
  f(x) = \pi^\nu x^{p^n} + b_{n-1}x^{p^{n-1}} + \dots + b_1 x^p + b_0 x
\]
\[
  f : \mathbb{G}_{a,S} \longrightarrow \mathbb{G}_{a,S}
\]
\[
  H_\eta = \operatorname{Ker} f_\eta .
\]
Par ailleurs $f$ est \struck{\ill{}} fid.\ plat --- car on \ill{} sur les fibres :
\uncertain{fibre par fibre}, et sur la fibre spéciale \uncertain{ceci provient} du fait
que $f_s$ n'est pas nul (car un des coeff.\ $b_{i s}$ est $\neq 0$). Donc
$\operatorname{Ker} f$ est un sous-schéma en groupes \add{fini} plat de $\mathbb{G}_{a,S}$,
de fibre générique $H_\eta$, donc coïncide avec $H$
\struck{Donc} (qui \uncertain{est} l'adhérence schématique de $H_\eta$ dans $\mathbb{G}_{a,S}$).
Donc $\operatorname{Ker} f_s$ est de rang $p^n$, or il est de rang $p^{n-i}$ si
$b_{i,s}$ est le premier des $b_{0,s}, b_{1,s}, \dots$ qui soit non nul.
C'est absurde.

\textbf{Lemme.} Soit $X_n^{\text{ét}} \subset X_n$ l'ouvert correspondant
\marginal{à justifier}
aux sous-schémas en groupes étales de $\mathbb{G}_a$. Alors $\varphi_n$ induit iso
\[
  \mathbb{E}^* \times \mathbb{E}^{n-1} \xrightarrow{\ \sim\ } X_n^{\text{ét}} .
\]

\page{5}
(on identifie $\mathbb{E}^* \times \mathbb{E}^{n-1}$ à l'ouvert de $\mathbb{E}^n$ dont
les \struck{\ill{}} $k$-points sont les $(a_0, \dots, a_{n-1}) \in k^n$ tels que
$a_0 \in k^*$ \ill{}).

\textbf{Question.} $X_n$ est-il réduit (i.e. \uncertain{immersion} $\subset \mathbb{E}^n$) ?
Ou en \uncertain{contient} l'un de ses points réduits \ill{} il réduit :
$X_n^{\mathrm{réd}} \simeq \mathbb{E}^* \times \mathbb{E}^{n-1}$ ?

Cas $n = 1$, $X_1$ est-il réduit \add{i.e. lisse} \uncertain{au point $0$} ?
Il revient au \uncertain{même} de dire que l'espace tangent de Zariski de $X_1$
en $0$ est de dim $1$ (à priori c'est de dim $\geqslant 1$). \struck{Donc} Il faut
étudier les points de $X_0$ \uncertain{à valeurs dans} $k = \mathbb{F}_p[J]/J^2$
se réduisant suivant $0$, i.e. les sous-groupes plats \uncertain{$H$} de
$\mathbb{G}_{a,k}$ tels que $H_0 = \alpha_p$.
\note{$k$ est l'algèbre des nombres duaux ; la lecture de la barre de fraction est douteuse}

\[
  \begin{array}{lll}
    H \subset G & H_0 \subset G_0 \to K_0 \simeq G_0/H_0 & \\
    G_0 & \overline{H_0} \subset G & \overline{K_0} \simeq G/\overline{H_0} \simeq \overline{G_0}/\overline{H_0}
  \end{array}
\]
\note{entre $G_0$ et $K_0$ de la première ligne, un mot biffé et illisible}
\[
  K = G/H \longrightarrow \text{une part.\ de } K_0 = G_0/H_0
\]
défini par une extension de $K_0$ par $J\,K_0$ (les splittings correspondant aux
\ill{} de $\overline{K_0} = G/\overline{H_0}$, $K_0 = G/H$ \struck{\ill{}}) qui partagent
\ill{} l'heure identique)
\marginal{\ill{} ; Zariski et \ill{} ; \uncertain{plat de \ill{}} ; \ill{}}
\ill{} en \uncertain{raison} de cette extension par $G \to K_0$
est l'obstruction : partager $G \to K_0$
en $G \to \overline{K_0}$, ou aller si se partage des
\uncertain{façons canoniques}, \ill{} $G$ dans cette extension

\page{6}
\ill{} \uncertain{splitting}, d'où un \uncertain{hom.} de $H_0$ dans $J K_0$ qui
donne l'extension de $K_0$ par $J K_0$ plus l'\uncertain{hom.} $G \to K$ \ill{}
(donc \uncertain{provenant} de \ill{} $H$)
\[
  u : H_0 \longrightarrow J K_0 .
\]
\add{Ici} $H_0 = \alpha_p$ $(\subset G_0 = \mathbb{G}_a)$, $K_0 \simeq \mathbb{G}_a$,
$J K_0 \simeq \mathbb{G}_a$, $\operatorname{Hom}(\alpha_p, \mathbb{G}_a)$
et \ill{} sont \uncertain{$\operatorname{Hom}(\alpha_p, \text{faisceau})
\xrightarrow{\ \sim\ }$} \uncertain{$\operatorname{Hom}(J\alpha_p, J\,\text{faisceau})$}
donc $\operatorname{Hom}(H_0, J K_0) \to \operatorname{Hom}(H_0, J K_0)$ \uncertain{iso},
donc \struck{l'extension} \uncertain{$\overline{K}$} de $K_0$ est triviale.
On trouve l'espace tangent de Zariski
\[
  : \quad \operatorname{Hom}_{\mathbb{F}_p}(H_0, J K_0) .
\]

\section*{Théorème : $\mathbb{E}^n \xrightarrow{\ \sim\ } X_n$}

\textbf{Théorème.} $\varphi_n : \mathbb{E}^n_{\mathbb{F}_p} \longrightarrow X_n$
\emph{iso}

i.e. si $S$ est un schéma de car.\ $p$, $H \subset \mathbb{G}_{a,S}$
sous-schéma en groupes fini et plat \struck{\ill{}}
\marginal{\uncertain{localement libre} de rang $p^n$}
alors
\[
  \exists\,! \ \ a = (a_0, \dots, a_{n-1}) \in \mathbb{E}^n(S) = \Gamma(S, \mathcal{O})^n
\]
tel que
\[
  H = \operatorname{Ker} f_a
\]
où
\[
  f_a : \mathbb{G}_{a,S} \to \mathbb{G}_{a,S}, \qquad
  f_a(x) = x^{p^n} + a_{n-1}x^{p^{n-1}} + \dots + a_0 x .
\]
De plus, $H$ est étale $/S$ sss $a_0 \in \Gamma(S, \mathcal{O})^* = \Gamma(S, \mathcal{O}^*)$.

\section*{Le déterminant $\Delta_n$ et les invariants de $GL(n,\mathbb{F}_p)$}

\page{7}
$\mathbb{E}^n$ sur $\mathbb{F}_p$, $\Gamma_n = GL(n, \mathbb{F}_p)_{\mathbb{F}_p}
\subset GL(n)_{\mathbb{F}_p}$
\marginal{opère sur $\mathbb{E}^n$}
\[
  Q_n = \mathbb{E}^n/\Gamma_n = \operatorname{Spec}(A_n), \qquad
  A_n = \mathbb{F}_p[X_1, \dots, X_n]^{\Gamma_n} .
\]
Soit $U_n \subset \mathbb{E}^n$ l'ouvert formé des syst.\ « linéairement indépendants
sur $\mathbb{F}_p$ », i.e. le \struck{plus grand} ouvert de $\mathbb{E}^n$ où le groupe
$\Gamma_n$ opère \uncertain{librement}, $Q_n^0 = U_n/\Gamma_n$.

\textbf{Lemme.} On a \struck{\ill{}} dans $\mathbb{F}_p[X_1, \dots, X_n]$
\begin{align*}
  \det\begin{pmatrix}
    X_1 & \cdots & X_n \\
    X_1^{p} & \cdots & X_n^{p} \\
    \vdots & & \vdots \\
    X_1^{p^{n-1}} & \cdots & X_n^{p^{n-1}}
  \end{pmatrix}
  &= \prod_{(\alpha_1, \dots, \alpha_{n-1}) \in \mathbb{F}_p^{n-1}}\Bigl(X_n + \sum_1^{n-1}\alpha_i X_i\Bigr) \\
  &\qquad \prod_{(\alpha_1, \dots, \alpha_{n-2}) \in \mathbb{F}_p^{n-2}}\Bigl(X_{n-1} + \sum_1^{n-2}\alpha_i X_i\Bigr)\cdots
\end{align*}
\[
  \cdots \Bigl(\prod_{\alpha_1 \in \mathbb{F}_p}(X_2 + \alpha_1 X_1)\Bigr) X_1
  \;=\; \prod_{0 \leqslant j \leqslant n-1}\ \prod_{\alpha \in k^j}\Bigl(X_{j+1} + \sum_{i=1}^{j}\alpha_i X_i\Bigr)
\]
\marginal{noté $\Delta_n(X_1, \dots, X_n)$}
\note{le second produit est indexé par $\alpha \in k^j$ sur le feuillet, là où le
premier membre est écrit sur $\mathbb{F}_p$}

\textbf{N.B.} À un facteur $\in \mathbb{F}_p^*$ près, le déterminant est le produit
des combinaisons linéaires non \struck{\ill{}} id.\ nulles des $X_i$ à coeff.\ dans
$\mathbb{F}_p$, \struck{\ill{}} ces combinaisons, modulo \ill{}
\struck{\ill{}} \uncertain{multiplication} par un $c \in \mathbb{F}_p^*$.
C'est un polynôme de degré
\[
  1 + p + \dots + p^{n-1} = \frac{p^n - 1}{p-1} = \operatorname{card} \mathbb{P}^n(\mathbb{F}_p)
\]
\note{le cardinal de $\mathbb{P}^{n-1}(\mathbb{F}_p)$ est ce que donne la formule ; il
écrit $\mathbb{P}^n$}
donc de degré égal \uncertain{au précédent}, $Q$. En tous cas, les zéros
\struck{des produits} \add{de chaque facteur $P_\lambda$ de $P$} sont contenus dans
celui de $Q$, $(P_\lambda$ étant premier$)$ \struck{\ill{}} donc $P_\lambda$ divise
$Q$, donc (les $P_\lambda$ étant premiers $2$ à $2$ distincts \uncertain{non associés})
$Q$ est divisible par $P = \prod_\lambda P_\lambda$, donc, étant de même degré, on a
$Q = cP$, $c \in \mathbb{F}_p^*$. Reste à déterminer $c$.

\page{8}
\textbf{Corollaire.} Soit $k$ corps de car.\ $p$,
$\xi = (\xi_1, \dots, \xi_n) \in k^n$. Alors
$\Delta_n(\xi_1, \dots, \xi_n) = 0$ sss $\xi_1, \dots, \xi_n$ sont lin.\ dép.\ sur
$\mathbb{F}_p$.

\textbf{Corollaire.} $U_n = \mathbb{E}^n_{\Delta_n} = \operatorname{Spec}
k[X_1, \dots, X_n]_{\Delta_n}$.

Si $\xi_1, \dots, \xi_n \in k$ corps de car.\ $p$, cherchons
$a = (a_0, \dots, a_{n-1}) \in k^n$ tel que l'on ait
\[
  (*) \qquad \xi_i^{p^n} + a_{n-1}\xi_i^{p^{n-1}} + \dots + a_1 \xi_i^{p} + a_0 \xi_i = 0
\]
i.e.
\[
  f_a(\xi_i) = 0, \qquad 1 \leqslant i \leqslant n .
\]
C'est un système de $n$ équations linéaires \struck{\ill{}} \add{en les $a_i$} aux
seconds membres $-\xi_i^{p^n}$, dont le déterminant est
$\Delta_n(\xi_1, \dots, \xi_n)$. Sous réserve que celui-ci est inversible
(i.e. les $\xi_i$ engendrent un sous-groupe \add{fini} (étale) de $\mathbb{G}_{a,k}$
de rang $p^n$) il y a \uncertain{une solution et une seule}, donnée par les formules
de Cramer
\[
  a_0 = \frac{1}{\Delta_n(\xi)}\det\begin{pmatrix}
    \xi_1^{p^n} & \cdots & \xi_n^{p^n} \\
    \xi_1^{p} & \cdots & \xi_n^{p} \\
    \vdots & & \vdots \\
    \xi_1^{p^{n-1}} & \cdots & \xi_n^{p^{n-1}}
  \end{pmatrix}
  = \frac{\Delta_0(\xi)}{\Delta_n(\xi)}
\]
\[
  a_i = \frac{1}{\Delta_n(\xi)}\det\begin{pmatrix}
    \xi_1 & \cdots & \xi_n \\
    \vdots & & \vdots \\
    \xi_1^{p^{i-1}} & \cdots & \xi_n^{p^{i-1}} \\
    \xi_1^{p^{n}} & \cdots & \xi_n^{p^{n}} \\
    \xi_1^{p^{i+1}} & \cdots & \xi_n^{p^{i+1}} \\
    \vdots & & \vdots \\
    \xi_1^{p^{n-1}} & \cdots & \xi_n^{p^{n-1}}
  \end{pmatrix}
  = \frac{\Delta_i(\xi)}{\Delta_n(\xi)}
\]
\marginal{\uncertain{modulo les lignes}, dont \ill{} \uncertain{suit par suite}}
\[
  a_{n-1} = \frac{1}{\Delta_n(\xi)}\det\begin{pmatrix}
    \xi_1 & \cdots & \xi_n \\
    \vdots & & \vdots \\
    \xi_1^{p^{n-2}} & \cdots & \xi_n^{p^{n-2}} \\
    \xi_1^{p^{n}} & \cdots & \xi_n^{p^{n}}
  \end{pmatrix}
  = \frac{\Delta_{n-1}(\xi)}{\Delta_n(\xi)}
\]
\note{les lignes des mineurs sont en partie illisibles ; elles sont restituées ici
selon la règle de Cramer que la page énonce, la ligne $\xi_i^{p^{n}}$ venant
remplacer la ligne $\xi_i^{p^{j}}$ d'indice correspondant}

\page{9}
\textbf{Lemme.} Prenons $k = \mathbb{F}_p[X_1, \dots, X_n]$, $\xi_i = X_i$, donc les
\struck{\ill{}} $\Delta_i \in \mathbb{F}_p[X_1, \dots, X_n]$
$(0 \leqslant i \leqslant n)$ donc, \ill{}
\struck{Alors} \add{on a}
\[
  \Delta_n X_i^{p^n} + \Delta_{n-1} X_i^{p^{n-1}} + \dots + \Delta_1 X_i^{p} + \Delta_0 X_i = 0
  \qquad (1 \leqslant i \leqslant n) .
\]
Les $\Delta_i$ $(0 \leqslant i \leqslant n-1)$ sont divisibles par $\Delta_n$, soit
\[
  \Delta_i = A_i \Delta_n \qquad 0 \leqslant i \leqslant n-1
\]
et
\[
  \Delta_0 = \pm\,\Delta_n^{p}\,(-1)^n
\]
\uncertain{au signe près}
\note{la puissance de $\Delta_n$ est ici lue $p$ ; le corollaire qui suit
donne $\Delta_n^{p-1}$ pour $A_0$, ce qui est cohérent avec $\Delta_0 = A_0\Delta_n$}
\marginal{et \ill{} on a donc}

\textbf{Cor.\ 1.}
\[
  X_i^{p^n} + A_{n-1}X_i^{p^{n-1}} + \dots + A_1 X_i^{p} + A_0 X_i = 0
  \qquad (1 \leqslant i \leqslant n)
\]
et les $A_i \in k[X_1, \dots, X_n]$ $(0 \leqslant i \leqslant n-1)$ sont
caractérisés par \uncertain{cette} relation. \struck{En effet} $A_0 = \pm\,\Delta_n^{p-1}$
\struck{au signe près}

\textbf{Cor.\ 2.} Revenons à une \struck{\ill{}} \add{corps $k$ de} car.\ $\mathbb{F}_p$.
Pour \uncertain{que} l'on ait $\Delta_n(\xi) \neq 0$ i.e. $\xi_1, \dots, \xi_n$
lin.\ indép.\ sur $\mathbb{F}_p$, il faut et il suffit \uncertain{que} $\Delta_0(\xi) \neq 0$.
\marginal{Cor.\ 2. Les $A_i$ sont \ill{} \struck{\ill{}} des invariants de
$GL(n,\mathbb{F}_p) = \Gamma_n$. \uncertain{La factorisation} \ill{} ; $A_i$ est
\uncertain{homogène et de degré} \ill{}}

Considérons
\[
  \mathbb{E}^n \xrightarrow{\ A_*\ } \mathbb{E}^n
\]
défini par $A_0, \dots, A_{n-1}$. Il se factorise en \ill{}
\struck{\ill{}}
\[
  U_n \ \bigl(= \mathbb{E}^n_{\Delta_n}\bigr) = A_*^{-1}\bigl(\mathbb{E}^* \times \mathbb{E}^{n-1}\bigr)
\]
et $A_*$ induit \struck{\ill{}} iso

\begin{tikzcd}[column sep=small, row sep=small, nodes={font=\scriptsize}]
\mathbb{E}^n \arrow[rr, "A_*"] \arrow[rd, "\mathrm{can}"'] & & \mathbb{E}^n \\
 & Q_n \arrow[ru, "\varphi_n"'] &
\end{tikzcd}

\struck{et induit} Ceci dit.

\page{10}
Puis
\[
  \varphi_n^{-1}\bigl(\underbrace{\mathbb{E}^* \times \mathbb{E}^{n-1}}_{(a_0, \dots, a_{n-1}) \text{ tels que } a_0 \text{ inv.}}\bigr) = Q_n^0
\]
et
\[
  Q_n^0 \xrightarrow{\ \sim\ } \mathbb{E}^* \times \mathbb{E}^{n-1} .
\]
\struck{En effet} \supplied{[}Il est immédiat que $Q_n^0 \simeq$ « sous-groupes finis
étales de $\mathbb{G}_a$ de rang $p^n$ », ce qui donne une
\uncertain{interprétation} géométrique au foncteur \uncertain{constant}
$\mathbb{E}^* \times \mathbb{E}^{n-1}$, en \uncertain{termes de noyaux d'itérés}
$f_a$ \ill{}\add{]}

En effet, on a un \struck{foncteur} \uncertain{hom.} en sens inverse
\[
  Q_n^0 \longleftarrow \mathbb{E}^* \times \mathbb{E}^{n-1} \subset \mathbb{E}^n
\]
\[
  f_a \longleftarrow a
\]

\textbf{Thm.} $Q_n \longrightarrow \mathbb{E}^n$ est un isomorphisme,
ou d'autres formes
\[
  \mathbb{F}_p[X_1, \dots, X_n]^{\Gamma_n} \simeq \mathbb{F}_p[A_0, \dots, A_{n-1}] .
\]

Dès \uncertain{lors} déjà que $Q_n \to \mathbb{E}^n$ est birationnel,
il suffit de prouver qu'il est fini
\marginal{ou même propre}
\struck{\ill{}} $Q_n$ est intègre et $\mathbb{E}^n$ normal, en conclure
que c'est un iso. \struck{Il suffit de prouver}
\marginal{Ou aussi il suffit}
\struck{\ill{}} de prouver que $Q_n \to \mathbb{E}^n$ est quasi-fini (car on
\uncertain{a une immersion ouverte par le}

\page{11}
« Main Theorem » de Zariski) et surjectif ; et il suffit
de prouver \uncertain{la} $=$ chose pour $\mathbb{E}^n \xrightarrow{\ A_*\ } \mathbb{E}^n$,
donc prouver que pour $k$ corps algébriquement clos, l'application
\struck{\emph{Surjectivité}}
\[
  k^n \xrightarrow{\ A_*\ } k^n
\]
est \emph{surjective à fibres finies}.

1) \emph{Surjection.} Soit $a = (a_0, \dots, a_{n-1}) \in k^n$.
Si $a = 0$, c'est l'image de $0 = (\underbrace{0, \dots, 0}) \in k^n$.
\struck{\ill{}}
Sinon, soit $a_r$ le premier parmi $a_0, a_1, \dots$ qui soit $\neq 0$.
Considérons $a_r, \dots, a_{n-1} \in k^{n-r}$, il provient de
$\xi_1, \dots, \xi_{n-r}$ ; on complète par des $0$ pour avoir
$\xi_1, \dots, \xi_n$, ça marche, grâce au

\textbf{Lemme.}
\[
  \begin{cases}
    A_0(X_1, \dots, X_{n-1}, 0) = 0 \\
    A_i(X_1, \dots, X_{n-1}, 0) = A_{i-1}(X_1, \dots, X_{n-1})
  \end{cases}
  \qquad \text{pour } 1 \leqslant i \leqslant n
\]
donc
\[
  A_\alpha\bigl[X_1, \dots, X_{n-s}, \underbrace{0, \dots, 0}_{s}\bigr] =
  \begin{cases}
    0 & \text{si } 0 \leqslant \alpha < s \\
    A_{\alpha - s}(X_1, \dots, X_{n-s}) & \text{si } s \leqslant \alpha \leqslant n
  \end{cases}
\]

2) \emph{À fibres finies.} En fait, on voit tout de suite que les fibres
sont les trajectoires de $\Gamma_n$. [Ainsi, indépendamment du M.\ Th.\ de Zar., on
voit que $Q_n \to \mathbb{E}^n$ est surjectif, radiciel, birationnel.]

\textbf{N.B.} On peut aussi regarder les invariants de
$\mathbb{F}_p[X_1, \dots, X_n]$ sous $S\Gamma_n = SL(n, \mathbb{F}_p)$.
Parmi ces invariants il y a $\Delta_n$ (car \ill{} $\Delta$ \uncertain{en fait}

\page{12}
\[
  \Delta_n(u \cdot X) = (\det u)\,\Delta_n(X) \qquad ]
\]
\marginal{$u \in \Gamma_n$}
\struck{\ill{} coefficient}

\struck{donc les $\Delta_i = \Delta_n A_i$ sont aussi des invariants sous $S\Gamma_n$.}

\textbf{Corollaire.} \struck{\ill{}}
\[
  \mathbb{F}_p[X_1, \dots, X_n]^{S\Gamma_n} = \mathbb{F}_p[A_1, \dots, A_{n-1}, \Delta_n]
\]
$\bigl($et $\Gamma_n/S\Gamma_n \xrightarrow[\ \det\ ]{\ \sim\ } \mathbb{F}_p^*$ opère
sur $\mathbb{F}_p[A_1, \dots, A_{n-1}, \Delta_n]$ en invariant $A_1, \dots, A_{n-1}$
et par homothéties sur $\Delta_n\bigr)$.

\textbf{Dém.}
\marginal{\uncertain{plus $i^0$ iso}}
On a le diagramme commutatif (défini par l'inclusion
$\mathbb{F}_p[A_1, \dots, A_{n-1}, \Delta_n]$
$\subset \mathbb{F}_p[X_1, \dots, X_n]^{S\Gamma_n}$)

\begin{tikzcd}[column sep=small, row sep=small, nodes={font=\scriptsize}]
\mathbb{E}^n/S\Gamma_n \arrow[r, "i^0"] \arrow[d, "p = \mathrm{can}"'] & \mathbb{E}^n \arrow[d, "q"] \\
\mathbb{E}^n/\Gamma_n \arrow[r, "\sim"] & \mathbb{E}^n
\end{tikzcd}

$q$ défini par l'inclusion
$\mathbb{F}_p[A_1, \dots, A_{n-1}, \Delta_n^{p-1}] \subset \mathbb{F}_p[A_1, \dots, A_{n-1}, \Delta_n]$.
Ainsi $p$ et $q$ sont des \uncertain{revêtements} de groupe $\mathbb{F}_p^*$,
identifiant le but au quotient de la source par ledit groupe, et $i^0$ est une
\ill{} compatible avec ces opérations, donc génériquement bijectif et
birationnel, donc (comme $\mathbb{E}^n$ est normal) un iso par le Main Theorem.

\section*{Structures linéaires, racines $(p-1)$-ièmes, et la version « faisceau »}

\page{13}
\textbf{Corollaire.} Soit $M \subset \mathbb{G}_{a,S}$ un sous-groupe fini et plat
\struck{loc.\ libre} de rang $p^n$, d'où $a = (a_0, \dots, a_{n-1}) \in \Gamma(S, \mathcal{O})^n$
tel que $M = \operatorname{Ker} f_a$, \add{$a_0$ inversible.}
On a alors une correspondance
\marginal{précise}
biunivoque entre l'ens.\ des \emph{structures spéciales linéaires}
(sur $\mathbb{F}_p$) de $M$, et l'ens.\ des racines $(p-1)$-ièmes de $a_0$.
On \uncertain{aura} : $\bigl(\overset{n}{\Lambda}_{\mathbb{F}_p} M\bigr)^*$ s'identifie
à : $(\mathbb{F}_p^*)_S$-torseur des racines $(p-1)$-ièmes de $a_0$,
i.e. \struck{à} \uncertain{l'image inverse} de $a_0 \in \Gamma(S, \mathcal{O})^*$
par $\mathbb{G}_m \xrightarrow{\ x \mapsto x^{p-1}\ } \mathbb{G}_m$, qui est un
torseur sous $\mu_{p-1,S} \simeq (\mathbb{F}_p^*)_S$.

\textbf{N.B.} \struck{\ill{}} La donnée d'un \uncertain{sous-groupe} fini \ill{} $M$
de rang $n$ sur $\mathbb{F}_p$, sur $S$, équivaut : la donnée de $E$ loc.\ libre de
rang $n$, avec iso
\[
  E'^{(p)} \xrightarrow[\ \sim\ ]{\ u\ } E .
\]
(On prend $E \simeq M \otimes_{\mathbb{F}_p} \mathcal{O}_S$, \uncertain{on retrouve} $M$
comme $M = \operatorname{Ker}\bigl(V(E) \xrightarrow{\ \operatorname{id} - u \circ F_E\ } V(E)\bigr)$.)
La donnée d'un \uncertain{hom.} $M \to \mathbb{G}_a$ équivaut à la donnée d'un
\uncertain{hom.} $\mathcal{O}$-lin.\ $E \to \mathcal{O}_S$, i.e. d'une section $\xi$
de $\check{E}$. Notons que $\check{E}$ est aussi muni d'une structure
$\check{E}^{(p)} \xrightarrow{\ \check{u}\ } \check{E}$.

\page{14}
Pour \uncertain{une} section $\xi$ de $\check{E}$ (lié à un hom.\ $\varphi : M \to \mathbb{G}_a$)
on peut considérer la section
\[
  \xi^{(p)} \overset{\text{df}}{=} u\bigl(\xi^{(F)}\bigr)
\]
\marginal{\uncertain{par linéarité}}
et par itération les $\xi^{(p^i)}$. \struck{Telles} \uncertain{Elles} correspondent
aux composés de $\varphi : M \to \mathbb{G}_a$ avec
$F_{\mathbb{G}_a} : \mathbb{G}_a \to \mathbb{G}_a$, et \add{les} itérés
$F^i_{\mathbb{G}_a}$.

Considérons
\[
  \xi \wedge \xi^{(p)} \wedge \dots \wedge \xi^{(p^{n-1})} \in \Gamma\bigl(\overset{n}{\Lambda}\check{E}\bigr) .
\]
On trouve ainsi \struck{une \uncertain{application} polynomiale}
\marginal{au-dessous de \ill{}}
\[
  \Delta^E_n : \check{E} \longrightarrow \overset{n}{\Lambda}\check{E}
\]
homogène de degré $1 + p + \dots + p^{n-1} = \dfrac{p^n - 1}{p - 1}$, d'où
\struck{\ill{}}
\[
  \Delta^E_n : V(E) \longrightarrow V\bigl(\overset{n}{\Lambda}E\bigr)
\]
\note{sous les deux membres il écrit $W(\check{E})$ et $W(\overset{n}{\Lambda}\check{E})$}
\struck{Pour} \struck{\ill{} degré \ill{}}
Considérons l'image inverse dans $V(E)$ de la section nulle de
$V(\overset{n}{\Lambda}E)$ ; \struck{\ill{}} c'est un diviseur de Cartier relatif,
dont le complémentaire $V(E)^0$ est un \uncertain{ouvert} schém.\ dense.
Si $\xi$ est une section de cet \uncertain{ouvert}, les
$\xi, \xi^{(p)}, \dots, \xi^{(p^{n-1})}$ forment une base de $\check{E}$, et on peut
écrire
\[
  \xi^{(p^n)} + A_{n-1}(\xi)\,\xi^{(p^{n-1})} + \dots + A_0(\xi)\,\xi = 0
\]

\page{15}
avec des $A^E_i(\xi)$ $(0 \leqslant i \leqslant n-1)$ uniquement
\uncertain{déterminés}. On trouve ainsi des hom.
\[
  A^E_i \in \Gamma\bigl(V(E)^0, \mathcal{O}_{V(E)^0}\bigr)
  \quad \text{i.e.} \quad A^E_i : V(E)^0 \longrightarrow \mathbb{G}_{a,S}
\]
uniquement déterminés par les conditions précédentes.
On peut les interpréter comme des éléments de
\[
  \underline{\operatorname{Sym}}^\bullet(E)_{\Delta^E_n} .
\]
\struck{$A_i$ est homogène de degré $p^i$}

Ceci dit, on trouve que les $A_i$ sont des sections de
$\underline{\operatorname{Sym}}(E)$, i.e. sont définis sur $V(E)$
\uncertain{tout entier} ; bien entendu ils y sont
\uncertain{caractérisés} par les identités précédentes \struck{\ill{}}.
On a encore
\[
  A^E_0 = \bigl(\Delta^E_n\bigr)^{p-1}\,(\pm 1)
\]
\supplied{[}ceci a un sens, grâce à : $\bigl(\overset{n}{\Lambda}\check{E}\bigr)^{p-1} \simeq \mathcal{O}_S$,
i.e. $\bigl(\overset{n}{\Lambda}\check{E}\bigr)^{p} = \overset{n}{\Lambda}\check{E}'^{(p)}
\simeq \bigl(\overset{n}{\Lambda}\check{E}\bigr)$\add{]}

\textbf{N.B.} \struck{\ill{}} Supposons que $\xi$ soit une section de $V(E)^0$.
\struck{Alors} \add{Alors} les $\xi^{(p^i)}$ $(0 \leqslant i \leqslant n-1)$
forment une base de $\check{E}$, qui s'identifie à $\mathcal{O}_S^n$,
\uncertain{avec le choix de} \ill{} base \ill{}
\[
  e_0^{(p)} = e_1, \quad e_1^{(p)} = e_2, \quad \dots, \quad e_{n-2}^{(p)} = e_{n-1}, \quad \text{et}
\]

\page{16}
\[
  e_{n-1}^{(p)} = -\Bigl(\underbrace{A_0(\xi)}_{a_0}\,e_0 + \dots + \underbrace{A_{n-1}(\xi)}_{a_{n-1}}\,e_{n-1}\Bigr)
\]
i.e. on trouve la matrice
\[
  \begin{pmatrix}
    0 & 0 & \cdots & 0 & -a_0 \\
    1 & 0 & \cdots & 0 & -a_1 \\
    0 & 1 & & \vdots & \vdots \\
    \vdots & & \ddots & 0 & \\
    0 & 0 & \cdots & 1 & -a_{n-1}
  \end{pmatrix}
\]
\note{la dernière colonne est en partie recouverte de ratures ; seuls $-a_0$, $a_1$
et $-a_{n-1}$ se lisent avec certitude}
\uncertain{par} l'\uncertain{application} \ill{} précédente. Ainsi \struck{\ill{}}
\struck{\ill{}} \uncertain{retrouve-t-on que} le choix d'un \ill{}

Il dresse ensuite un tableau à deux colonnes, chacune de trois étages reliés par des
flèches montantes, les deux colonnes étant reliées par des inclusions
$\subset$ de gauche à droite.

\emph{Colonne de gauche} --- en haut : $(a_0, \dots, a_{n-1})$, $a_0$ inv. ; au
milieu : sous-groupes finis étales de $\mathbb{G}_a$ ; en bas : $E$ loc.\ libre $/S$
avec $\check{E}^{(p)} \simeq \check{E}$, et section $\xi$ de $\check{E}$ telle que
$\xi, \xi^{(p)}, \dots, \xi^{(p^{n-1})}$ soit une base.

\emph{Colonne de droite} --- en haut : $(a_0, \dots, a_n)$ ; au milieu :
sous-groupes finis loc.\ libres de rg $p^n$ de $\mathbb{G}_a$ ; en bas : $F$ loc.\
libre de rg $n$ avec $F^{(p)} \to F$ (pas nécess.\ iso) et $\xi \in \Gamma F$ telle
que $\xi, \xi^{(p)}, \dots, \xi^{(p^{n-1})}$ base.

\note{la flèche montante du bas vers le milieu, dans la colonne de droite, est
suivie d'un « ? » : c'est l'implication qu'il ne tranche pas. Une accolade court
du bas de la colonne de droite vers le haut, annotée \uncertain{« trivial ? »}}

\section*{Descente le long d'un pincement}

\page{18}

\begin{tikzcd}
X' \arrow[d] \\ X
\end{tikzcd}

\marginal{\uncertain{déduit} de $X'$ en « pinçant » $a$ et $b$}

\[
  H \subset \mathbb{G}_{a,X} = \operatorname{Ker}\Bigl(\mathbb{G}_{a,X}
  \xrightarrow{\ x^{p^n} + a_{n-1}x^{p^{n-1}} + \dots + a_0 x\ } \mathbb{G}_{a,X}\Bigr)
\]
\[
  H' \subset \mathbb{G}_{a,X'}, \qquad H'_a = H'_b, \qquad c_i(a) = c_i(b)
\]
\[
  \cdot \longrightarrow H \longrightarrow \mathbb{G}_{a,X} \longrightarrow \mathbb{G}_{a,X} \longrightarrow \cdot
\]
$\varphi|\hat{H} = 0$ i.e. $\hat{\varphi}$ se \struck{\ill{}} factorise en $\psi$
\note{deux flèches partent de ce point : l'une, notée $q$, descend vers
$\mathbb{G}_{a,X}$ ; l'autre, en pointillé, notée $\hat{r}$, y remonte}
\[
  (\operatorname{Ker}\varphi')_{y'} = (\operatorname{Ker}\varphi)_y = N'_y,
  \qquad \overline{N_y} = N'
\]
\marginal{adhérence schématique dans $\mathbb{G}_{a,X'}$}
\[
  N' \subset H', \qquad \hat{N'} \simeq \hat{H'} \quad\text{i.e.}\quad
  \mathbb{G}_{a,X'}/N' \longrightarrow \mathbb{G}_{a,X'} \ \ \text{ét.}
\]
\emph{Mais} supposons $N'(a) \neq N'(b)$.

\[
  N' \subset H' \subset \mathbb{G}_{a,X'}
\]
\[
  H'(a) = H'(b), \qquad N'(a) \neq N'(b)
\]
\[
  \hat{N'} \simeq \hat{H'} \quad\text{i.e.}\quad
  \mathbb{G}_{a,X'}/N' \longrightarrow \mathbb{G}_{a,X'}/H' \ \ \text{ét.}
\]
\[
  \cdot \to N' \to \mathbb{G}_{a,X'} \xrightarrow{\ f_{\alpha'}\ } \mathbb{G}_{a,X'} \to \cdot
  \qquad \xrightarrow{\ u_\lambda\ } \mathbb{G}_{a,X}
\]
\note{$u_\lambda$ est annoté \uncertain{« hom.\ Verschiebung »}}
\[
  \cdot \to K' \to \mathbb{G}_{a,X'} \xrightarrow{\ f_{\beta'}\ } \mathbb{G}_{a,X'} \to \cdot
\]
\note{une flèche $u_{\lambda'}$ part de cette seconde ligne, et une flèche en
pointillé rejoint $\psi$ ; le point d'arrivée est noyé sous une tache d'encre}
\[
  u_{\lambda'} f_{\alpha'} = \lambda' f_{\alpha'}
\]
\struck{\ill{}} \uncertain{pour} $i : X$ \struck{\ill{}}
\note{la fin de la page est une esquisse : les deux suites exactes et les flèches
qui les relient sont tracées, mais l'argument qui devait les comparer n'est pas
écrit}

\section*{Fragments au crayon}

\page{19}
\note{les deux feuillets qui suivent sont au crayon, d'un tracé pâle et sans
rédaction suivie : ce sont des brouillons de calcul. Ce qui se lit est donné
ici ; le reste est laissé \ill{}}

\struck{\ill{}}
\[
  X^{p^n} + A_1(X_1, \dots, X_n)\,X^{p^{n-1}} + A_2(X_1, \dots, X_n)\,X^{p^{n-2}}
  + \dots + A_{n-1}(X_1, \dots, X_n)\,X
\]
\note{ici les $A_i$ sont indexés en sens inverse de celui des pages 9 à 15}

Un bloc écrit dans le sens de la hauteur occupe la marge gauche ; on y lit
$H_0$, $\mathbb{G}_0$, $K_0 = \mathbb{G}_0/H_0$, $G/H_0 = \overline{G_0}/\overline{H_0} = \overline{K_0}$,
$G/H = K$, $\mathbb{G}_0 \to K_0 \Rightarrow K_0$, puis
$\operatorname{Hom}(\alpha_{p^r}, \mathbb{G}_a) = \operatorname{Hom}(\alpha_p, \ldots)$,
dont le second argument est \ill{}, puis $V_F$ suivi d'un terme \ill{},
$F^r = 0$, \uncertain{$\mathbb{F}_p(\mathbb{F})/F^r$}.
Le reste du bloc est \ill{}.
\note{ce bloc reprend, au crayon, le calcul d'espace tangent des pages 5 et 6}

\[
  X_1 e_1 + \dots + X_n e_n, \qquad X_1^{p} e_1, \ \dots, \qquad X_1^{p^{n-1}} e_1, \ \dots
\]
\[
  X_1^{p^n} + a_1 X_1^{p^{n-1}} + a_2 X_1^{p^{n-2}} + \dots + a_n X_1 = 0
\]
\[
  f^{p^n} + a_1 f^{p^{n-1}} + \dots + a_{n-1} f^{p} + a_n f = 0
\]
\[
  f, \ f^{p}, \ \dots, \ f^{p^{n-1}} \qquad\text{en regard de}\qquad e_1, \ e_2, \ \dots, \ e_n
\]
\[
  e_1^{(p)} = e_2, \quad e_2^{(p)} = e_3, \quad \dots, \quad e_{n-1}^{(p)} = e_n,
  \quad e_n^{(p)} = -\bigl(a_1 e_n + a_2 e_{n-1} + \dots + a_n e_1\bigr)
\]
\marginal{$a_n \neq 0$}
\[
  M \ \text{loc.\ lib.\ t.f.}, \qquad M^{(p)} \xrightarrow{\ \sim\ } M, \qquad
  M \longrightarrow 0
\]
\[
  e_1, \dots, e_n \ ; \qquad e_i \longmapsto \sum_j \lambda_{ij} e_j
\]
\[
  M \longrightarrow M^{(p)} \longrightarrow M
\]

\page{20}
\[
  \bigl(f^{p-1} - g^{p-1}\bigr) X^{p^2} - \bigl(f^{p^2} - g^{p^2}\bigr) X^{p}
  + \bigl(f^{p^2} g^{p-1} - g^{p^2} f^{p-1}\bigr) X
\]
\note{les exposants de cette ligne sont les plus incertains du feuillet ; il
souligne les trois coefficients par des accolades}
\[
  \begin{cases}
    f^{p^3} + a f^{p^2} + b f^{p} + c f = 0 \\
    g^{p^3} + a g^{p^2} + b g^{p} + c g = 0 \\
    h^{p^3} + a h^{p^2} + b h^{p} + c h = 0
  \end{cases}
\]
\[
  \det\begin{pmatrix}
    X_1^{p^{n-1}} & \cdots & X_1^{p} & X_1 \\
    \vdots & & \vdots & \vdots \\
    X_n^{p^{n-1}} & \cdots & X_n^{p} & X_n
  \end{pmatrix}
\]
\marginal{$1 + p + \dots + p^{n-1} = \dfrac{p^n - 1}{p-1}$}
\struck{\ill{}}
\[
  \prod_{\mathbb{P}^n(\mathbb{F}_p)} \bigl(\lambda_1 X_1 + \dots + \lambda_n X_n\bigr),
  \qquad (\lambda_1, \dots, \lambda_n)
\]
\[
  \prod_{i=1}^{n} \Biggl(\ \prod_{\lambda \in \mathbb{F}_p^{n-i}}
  \bigl(X_i + \lambda_1 X_{i+1} + \dots + \lambda_{n-i} X_n\bigr)\Biggr)
\]
\[
  X^{p^2} Y^{p} - Y^{p^2} X^{p}, \qquad X^{p^2} Y - Y^{p^2} X
\]
\[
  \frac{\bigl(X^{p}Y - Y^{p}X\bigr)^{p}}{X^{p}Y - Y^{p}X} = \bigl(X^{p}Y - Y^{p}X\bigr)^{p-1}
\]
\marginal{$(X \wedge Y)$}
\[
  \left(\frac{X}{Y}\right)^{p^2} = \frac{X}{Y}
\]
\ill{}

\end{document}
